% Problem-section file following ../_template.tex.
\section{Lockability of two-way distillable entanglement}
\label{sec:problem-21}

\paragraph{Problem.}
Can discarding one local qubit reduce two-way distillable entanglement by an
arbitrarily large amount?  Let $D_{\leftrightarrow}(A:B)_\rho$ denote the
asymptotic singlet-distillation rate of $\rho_{AB}$ under local operations and
unrestricted two-way classical communication.  The question is whether there
are finite-dimensional states $\rho^{(r)}_{A_ra:B_r}$, with $\dim a=2$, such
that
\begin{equation}
  D_{\leftrightarrow}(A_ra:B_r)_{\rho^{(r)}}
  -D_{\leftrightarrow}(A_r:B_r)_{\operatorname{Tr}_a\rho^{(r)}}
  \xrightarrow[r\to\infty]{}\infty.
  \label{eq:p21-two-way-locking}
\end{equation}
Equation~\eqref{eq:p21-two-way-locking} requires an unbounded loss while the
discarded subsystem has fixed dimension two.

\subsection*{Status}

Unsolved

\subsection*{Source}

Horodecki et al. prove lockability of one-way distillable entanglement and
explicitly leave the two-way quantity unresolved
\sourcecite{ref:p21-locking}{HHHO05}.

\subsection*{Progress}

\begin{itemize}
  \item Horodecki, Horodecki, Horodecki, and Oppenheim constructed families
  in which discarding one qubit causes an unbounded loss of entanglement of
  formation, entanglement cost, logarithmic negativity, or one-way
  distillable entanglement.  These results do not establish
  Eq.~\eqref{eq:p21-two-way-locking}, which allows unrestricted two-way
  communication \sourcecite{ref:p21-locking}{HHHO05}.

  \item The same work proved that discarding one qubit decreases the relative
  entropy of entanglement by at most two ebits, but explicitly left
  lockability of unrestricted distillable entanglement open.  No corresponding
  universal bound is known for $D_{\leftrightarrow}$
  \sourcecite{ref:p21-locking}{HHHO05}.
\end{itemize}

\subsection*{References}

\begin{enumerate}[label={},leftmargin=4.8em,labelsep=0.5em]
  \footnotesize
  \item[\textup{[HHHO05]}]\label{ref:p21-locking}
  K. Horodecki, M. Horodecki, P. Horodecki, and J. Oppenheim,
  ``Locking Entanglement with a Single Qubit,''
  \emph{Physical Review Letters} \textbf{94}, 200501 (2005).
  \href{https://doi.org/10.1103/PhysRevLett.94.200501}{doi:10.1103/PhysRevLett.94.200501};
  \href{https://arxiv.org/abs/quant-ph/0404096}{arXiv:quant-ph/0404096}.
\end{enumerate}

\subsection*{Comment}

The remaining gap is to construct a family satisfying
Eq.~\eqref{eq:p21-two-way-locking} or prove a dimension-independent upper
bound on the loss of $D_{\leftrightarrow}$.  The one-way analogue is already
known to be lockable.

\subsection*{Tag}

Entanglement distillation; Entanglement measures;
Local operations and classical communication; Qubit systems

\subsection*{ID}

\texttt{op\_7ceb25c3fe5c9403}
